\documentclass[12pt, a4paper]{ctexart}
\usepackage[margin=2cm]{geometry}
\usepackage{fancyhdr}
\usepackage{amsmath,amssymb,amsthm}
\usepackage{enumitem}
\usepackage{graphicx}
\usepackage{tikz}
\usepackage{pgfplots}
\pgfplotsset{compat=1.18}

\setCJKmainfont{Hiragino Sans GB}
\setCJKsansfont{Heiti SC}

\pagestyle{fancy}
\fancyhf{}
\fancyhead[L]{高考数学满分突破·函数分析与高考}
\fancyhead[R]{测试卷}
\fancyfoot[C]{\thepage}
\renewcommand{\headrulewidth}{0.4pt}

\newcommand{\blank}[1]{\underline{\hspace{#1}}}

\begin{document}

\begin{center}
{\LARGE\bfseries 高考数学满分突破·函数分析与高考\quad 测试卷}\\[0.5em]
{\large 满分150分\quad 考试时间120分钟}
\end{center}

\vspace{1em}

\noindent\textbf{一、单项选择题（本大题共8小题，每小题5分，共40分。每小题只有一个选项符合题目要求。）}

\begin{enumerate}
\item 函数$f(x)=\ln(x-1)+\sqrt{3-x}$的定义域为\blank{3cm}。
\begin{enumerate}[label=\Alph*.]
\item $(1,3)$
\item $[1,3]$
\item $(1,3]$
\item $[1,3)$
\end{enumerate}

\item 函数$f(x)=x^2-2x+3$在区间$[0,3]$上的值域为\blank{3cm}。
\begin{enumerate}[label=\Alph*.]
\item $[2,6]$
\item $[3,6]$
\item $[2,3]$
\item $[2,+\infty)$
\end{enumerate}

\item 函数$f(x)=\dfrac{2^x-1}{2^x+1}$为\blank{3cm}。
\begin{enumerate}[label=\Alph*.]
\item 奇函数，在$\mathbb{R}$上单调递增
\item 偶函数，在$\mathbb{R}$上单调递增
\item 奇函数，在$\mathbb{R}$上单调递减
\item 偶函数，在$\mathbb{R}$上单调递减
\end{enumerate}

\item 已知$a=2^{0.3}$，$b=\log_2 0.3$，$c=0.3^2$，则$a,b,c$的大小关系为\blank{3cm}。
\begin{enumerate}[label=\Alph*.]
\item $a>b>c$
\item $a>c>b$
\item $c>a>b$
\item $b>a>c$
\end{enumerate}

\item 函数$f(x)=\ln x+x-2$的零点所在区间为\blank{3cm}。
\begin{enumerate}[label=\Alph*.]
\item $(0,1)$
\item $(1,2)$
\item $(2,3)$
\item $(3,4)$
\end{enumerate}

\item 设$f(x)$是定义在$\mathbb{R}$上的奇函数，当$x>0$时，$f(x)=x^2-2x$，则$f(-1)$的值为\blank{3cm}。
\begin{enumerate}[label=\Alph*.]
\item $-1$
\item $1$
\item $3$
\item $-3$
\end{enumerate}

\item 函数$f(x)=\dfrac{x}{x^2+1}$的最大值为\blank{3cm}。
\begin{enumerate}[label=\Alph*.]
\item $\dfrac{1}{2}$
\item $1$
\item $\dfrac{\sqrt{2}}{2}$
\item $2$
\end{enumerate}

\item 已知$f(x)$是定义在$\mathbb{R}$上的偶函数，且在$[0,+\infty)$上单调递减，则$f(x^2-2x)$的单调递增区间为\blank{3cm}。
\begin{enumerate}[label=\Alph*.]
\item $[0,1]$
\item $[1,+\infty)$
\item $(-\infty,0]$
\item $(-\infty,0]\cup[2,+\infty)$
\end{enumerate}
\end{enumerate}

\vspace{1em}

\noindent\textbf{二、多项选择题（本大题共4小题，每小题5分，共20分。每小题有多个选项符合题目要求，全部选对得5分，部分选对得2分，有选错的得0分。）}

\begin{enumerate}[start=9]
\item 关于函数$f(x)=\ln x$，下列结论正确的有\blank{3cm}。
\begin{enumerate}[label=\Alph*.]
\item $f(x)$在$(0,+\infty)$上单调递增
\item $f(x)$的值域为$\mathbb{R}$
\item $f(xy)=f(x)+f(y)$（$x,y>0$）
\item $f(x^n)=nf(x)$（$x>0$，$n\in\mathbb{R}$）
\end{enumerate}

\item 关于函数的奇偶性，下列说法正确的有\blank{3cm}。
\begin{enumerate}[label=\Alph*.]
\item 奇函数的图像关于原点对称
\item 偶函数的图像关于$y$轴对称
\item 若$f(x)$为奇函数且在$x=0$处有定义，则$f(0)=0$
\item $f(x)=x^2+1$既是奇函数又是偶函数
\end{enumerate}

\item 设$f(x)=\begin{cases} x^2, & x\leq 0 \\ \ln x, & x>0 \end{cases}$，则下列结论正确的有\blank{3cm}。
\begin{enumerate}[label=\Alph*.]
\item $f(x)$在$x=0$处连续
\item $f(x)$在$(-\infty,0)$上单调递减
\item $f(x)$在$(0,+\infty)$上单调递增
\item $f(x)$的值域为$[0,+\infty)$
\end{enumerate}

\item 关于指数函数$y=a^x$（$a>0,a\neq 1$），下列结论正确的有\blank{3cm}。
\begin{enumerate}[label=\Alph*.]
\item 图像恒过点$(0,1)$
\item 当$a>1$时，函数在$\mathbb{R}$上单调递增
\item 当$0<a<1$时，函数在$\mathbb{R}$上单调递减
\item 值域为$(0,+\infty)$
\end{enumerate}
\end{enumerate}

\vspace{1em}

\noindent\textbf{三、填空题（本大题共4小题，每小题5分，共20分。）}

\begin{enumerate}[start=13]
\item 函数$f(x)=\log_2(x^2-2x-3)$的单调递增区间为\blank{3cm}。

\item 已知$f(2x+1)=x^2-2x$，则$f(x)=$ \blank{4cm}。

\item 方程$\log_2(x+3)+\log_2(1-x)=1$的解为\blank{2cm}。

\item 函数$f(x)=2^x+x-4$的零点为\blank{2cm}。
\end{enumerate}

\vspace{1em}

\noindent\textbf{四、解答题（本大题共6小题，共70分。解答应写出文字说明、证明过程或演算步骤。）}

\begin{enumerate}[start=17]
\item （10分）已知函数$f(x)=\dfrac{2x+1}{x-1}$。

（1）求$f(x)$的定义域和值域；

（2）证明：$f(x)$在$(1,+\infty)$上单调递减；

（3）求$f^{-1}(x)$（反函数）。

\item （10分）已知函数$f(x)=\ln x-ax+\dfrac{a}{x}$（$a\in\mathbb{R}$）。

（1）当$a=1$时，求$f(x)$的单调区间；

（2）若$f(x)$在$[1,+\infty)$上单调递减，求$a$的取值范围。

\item （12分）已知函数$f(x)=x^2-2ax+3$（$a\in\mathbb{R}$）。

（1）若$f(x)$在$[1,+\infty)$上单调递增，求$a$的取值范围；

（2）若$f(x)$在$[-1,1]$上的最小值为$g(a)$，求$g(a)$的表达式；

（3）在（2）的条件下，求$g(a)$的最大值。

\item （12分）已知函数$f(x)=a^x$（$a>0,a\neq 1$）。

（1）若$f(2)=4$，求$a$的值，并解不等式$f(x)>f(x+1)$；

（2）设$g(x)=a^x-a^{-x}$，判断$g(x)$的奇偶性并证明；

（3）在（2）的条件下，若$g(x)$在$\mathbb{R}$上单调递增，求$a$的取值范围。

\item （13分）已知函数$f(x)=\log_a(x+1)-\log_a(1-x)$（$a>0,a\neq 1$）。

（1）求$f(x)$的定义域；

（2）判断$f(x)$的奇偶性并证明；

（3）当$a>1$时，求$f(x)>0$的解集。

\item （13分）已知函数$f(x)=x+\dfrac{a}{x}$（$a>0$）。

（1）利用导数，求$f(x)$在$(0,+\infty)$上的单调区间；

（2）求$f(x)$在$(0,+\infty)$上的最小值；

（3）利用（2）的结论，证明：当$x>0$，$y>0$时，$\dfrac{x^2+y^2}{xy}\geq 2$；

（4）更一般地，证明：当$x_i>0$（$i=1,2,\ldots,n$）时，$\sum\limits_{i=1}^{n}x_i\cdot\sum\limits_{i=1}^{n}\dfrac{1}{x_i}\geq n^2$。
\end{enumerate}

\end{document}
